% Four narrative paragraphs, then contributions and the six-section roadmap.
% First-pass text is preserved in first_pass/. Final evaluation is pending.
\label{sec:introduction}

Underwater wireless sensor networks support repeated data collection using
battery-powered sensors and acoustic communication. A clustering scheme
selects cluster heads (CHs), assigns ordinary sensors to them, and forwards
collected information toward a sink. The CH decision determines not only
collection distances but also reception, aggregation, and forwarding costs.
As rounds continue, the resulting differences in consumption produce
different residual energies across the network. Repeated CH selection must
therefore be considered together with the distribution of communication
work, rather than only the geometry of a single round. Energy-aware
clustering and multi-hop collection are established approaches to this
task~\cite{Xiao2020ACOCR,Nguyen2021EECMR}.

Residual energy, spatial distribution, and relay burden have already been
used in CH selection and unequal clustering~\cite{Hou2022UCPSO,Zhu2021SoftKMeans}.
The issue examined here is more specific: total modeled energy does not
uniquely describe the severity of the load assigned to individual nodes.
For example, spending $0.02$~J consumes one fifth of a node's remaining
$0.10$~J, but only one fiftieth of another node's remaining $1.00$~J.
Two complete routing trees with the same total consumption can thus
have different fractional-depletion patterns. Evaluating a completed tree
captures forwarding costs, yet ranking that tree solely by its total
energy still does not distinguish these patterns. This difference remains
relevant even when all compared methods use the same routing procedure.

Uneven deployment adds a spatial dimension to that decision. Local
support can guide the placement of CHs, but the support still available
after some CHs are selected differs from the structural support before
selection. Moreover, a node's neighborhood does not by itself reveal the
traffic assigned by a completed routing tree. This motivates investigating
two complementary uses of density: remaining support during candidate
construction and structural density when weighting evaluated fractional
depletion. The research question is whether this combination improves
modeled operation duration and energy use at an acceptable computational
cost under fixed-position, three-dimensional simulation assumptions.

To investigate this question, we propose density-aware ant colony
optimization (DACO). Its construction heuristic combines residual energy,
node pheromone, and a grid-support measure reduced as CHs are added.
Each candidate is decoded using the common member-assignment and
multi-hop routing procedure. Its score combines total consumption with
a powered fractional-depletion penalty divided by base structural
density. Reference scales are calibrated and held fixed within a round,
so that candidate comparisons use a common score definition. Completed-tree
evaluation and between-candidate node-pheromone deposits are shared with
the energy-only NodeACO reference implementation; the distinguishing
mechanisms studied here are density-guided construction and the calibrated
depletion-aware ranking rule. Section~4 specifies the remaining schedule
differences, while controlled comparisons are needed to identify their effects.

The contributions are organized around these mechanisms:
\begin{itemize}
\item A grid-support CH-construction heuristic that distinguishes base
structural density from remaining support after partial CH selection.
\item A calibrated candidate score combining total tree consumption and
density-weighted fractional depletion, with its interpretation and
comparison context explicitly specified.
\item \textbf{[Final evaluation required: replace this item with verified
findings from controlled density/score comparisons, deployment studies,
and computational-cost measurements.]}
\end{itemize}
% No claim that density-aware clustering or node pheromone is first proposed.

The rest of this paper is organized as follows. Section~2 reviews related
work. Section~3 presents the network model and problem formulation.
Section~4 describes DACO. Section~5 presents the experiments and their
analysis. Section~6 concludes the paper.
